\documentclass[../../main.tex]{subfiles}
\begin{document}

{\bf [解]}

\paragraph{(1)}
まず $a$ を固定する．$a > 0$ だから，$u=b/a,v=c/a$の動く範囲は
\begin{align}
  \begin{cases}
    \dfrac{2.7}{a} \le u \le \dfrac{3.3}{a} \\
    \dfrac{4.5}{a} \le v \le \dfrac{5.4}{a}
  \end{cases} \label{eq:1}
\end{align}
である．これを$(u,v)$平面に図示すると下図のような長方形$ABCD$の内部及び周になる．
各点の座標は
\begin{align}
  A\left(\dfrac{3.3}{a},\dfrac{4.5}{a}\right), \quad
  B\left(\dfrac{3.3}{a},\dfrac{5.4}{a}\right), \\
  C\left(\dfrac{2.7}{a},\dfrac{5.4}{a}\right), \quad
  D\left(\dfrac{2.7}{a},\dfrac{4.5}{a}\right),
\end{align}
である．

% TODO: 図

次に $0.9 \le a \le 1.1$ の範囲で $a$ を動かす．長方形の頂点$A, B, C, D$ はそれぞれ
\begin{align}
  \begin{cases}
    A: v = \dfrac{15}{11}u & \left(3\le u \le \dfrac{11}{3}\right) \\
    B: v = \dfrac{18}{11}u & \left(3\le u \le \dfrac{11}{3}\right) \\
    C: v = 2u              & \left(\dfrac{27}{11}\le u \le 3\right) \\
    D: v = \dfrac{5}{3}u   & \left(\dfrac{27}{11}\le u \le 3\right)
  \end{cases}
\end{align}
を動くので，動く範囲は下図斜線部 (境界含む)である．

\begin{figure}[htb]
  \centering
  \begin{tikzpicture}[scale=1.2, >=stealth]
    % Rectangle for fixed a
    \begin{scope}[shift={(-2.5,0)}, scale=0.9]
      \draw[thick] (0,0) rectangle (2,1.2);
      \node[below left] at (0,0) {$D(\frac{2.7}{a}, \frac{4.5}{a})$};
      \node[below right] at (2,0) {$A(\frac{3.3}{a}, \frac{4.5}{a})$};
      \node[above right] at (2,1.2) {$B(\frac{3.3}{a}, \frac{5.4}{a})$};
      \node[above left] at (0,1.2) {$C(\frac{2.7}{a}, \frac{5.4}{a})$};
    \end{scope}

    % Region in uv-plane
    \begin{scope}[shift={(2.5,0)}, scale=1.1]
      \draw[->] (-0.3,0) -- (3.8,0) node[right] {$u$};
      \draw[->] (0,-0.3) -- (0,3.2) node[above] {$v$};
      \node[below left] at (0,0) {$0$};

      % Points
      % u_1 = 27/11 = 2.45, u_2 = 3, u_3 = 11/3 = 3.67
      % v_1 = 45/11 = 4.09, v_2 = 2*u, etc.
      \draw[dashed] (2.45,0) node[below] {$\frac{27}{11}$} -- (2.45,2.45);
      \draw[dashed] (3.0,0) node[below] {$3$} -- (3.0,2.7);
      \draw[dashed] (3.67,0) node[below] {$\frac{11}{3}$} -- (3.67,2.5);

      % Shaded region
      \fill[pattern=north east lines, pattern color=gray] (2.45, 2.05) -- (3.0, 2.7) -- (3.67, 2.7) -- (3.67, 2.25) -- (3.0, 1.8) -- (2.45, 1.8) -- cycle;
      \draw[thick] (2.45, 2.05) -- (3.0, 2.7) -- (3.67, 2.7) -- (3.67, 2.25) -- (3.0, 1.8) -- (2.45, 1.8) -- cycle;

      \draw (0,0) -- (3.5, 2.87) node[right] {$v = \frac{15}{11}u$};
      \draw (0,0) -- (1.5, 3.0) node[above] {$v = 2u$};
    \end{scope}
  \end{tikzpicture}
  \caption{固定された $a$ に対する長方形 $ABCD$ と，$a$ を動かした時の $uv$ 平面上の存在領域}
\end{figure}

\paragraph{(2)}

まず $Z$ の存在条件から$ax^2-2bx+c=0$の判別式が 0 以上だから
\begin{align}
  b^2 - ac \ge 0
  \iff
  \Bigl(\dfrac{b}{a}\Bigr)^2 - \Bigl(\dfrac{c}{a}\Bigr) \ge 0 \quad (\because a > 0)
  \iff
  u^2 - v \ge 0 \label{eq:2}
\end{align}
である．

% TODO: (1)の領域+\cref{eq:2}の図

上図から，(1)で求めた領域を $(u, v)$ が動く時，\cref{eq:2}は常に満たされる．

そこで具体的に$z$を計算すると二次関数の解の公式から
\begin{align}
  z = \dfrac{b + \sqrt{b^2 - ac}}{a} = u + \sqrt{u^2 - v} \label{eq:3}
\end{align}
である．

ここで(1) で求めた領域は式にすると，
\begin{align}
  \begin{cases}
    \dfrac{27}{11} \le u \le 3 \text{ の時，} \dfrac{45}{11} \le v \le 2u \\
    3 \le u \le \dfrac{11}{3} \text{ の時，} \dfrac{15}{11}u \le v \le 6
  \end{cases} \label{eq:4}
\end{align}
である．$u$を固定したとき\cref{eq:3}で$z$は$v$の単調減少関数だから\cref{eq:4}の上端・下端でその最大・最小値をとる．すなわち
\begin{align}
  \begin{cases}
    \dfrac{27}{11} \le u \le 3 \text{ の時，} \min Z = u + \sqrt{u^2 - 2u}, \quad \max Z = u + \sqrt{u^2 - \dfrac{45}{11}} \\
    3 \le u \le \dfrac{11}{3} \text{ の時，} \min Z = u + \sqrt{u^2 - 6}, \quad \max Z = u + \sqrt{u^2 - \dfrac{15}{11}u}
  \end{cases} \label{eq:5}
\end{align}
である．

% TODO: z(u)のグラフ

次に $u$ を動かす．\cref{eq:5}はいずれも$u$の範囲では単調増加で，かつ$u=3$で連続であるから，$z$が最小値をとるのは$u=\dfrac{27}{11}$の時，最大値をとるのは$u=\dfrac{11}{3}$の時で，それぞれ値は
\begin{align}
  \begin{cases}
    \min z = z\left(\dfrac{27}{11}\right) = \dfrac{27+3\sqrt{15}}{11} \\
    \max z = z\left(\dfrac{11}{3}\right)  = \dfrac{11+2\sqrt{19}}{3}
  \end{cases}
\end{align}
である．

\end{document}
